\section{Sequences with repetition}

A sequence with repetition is an ordered list in which each position is filled independently from the same set of symbols.

\begin{definitionbox}
Let \(S\) be a set with \(n\) elements. A \textbf{sequence of length \(k\) with repetition} over \(S\) is an ordered \(k\)-tuple
\[
(a_1,a_2,\ldots,a_k),
\]
where each \(a_i\in S\). The same element may occur many times.
\end{definitionbox}

The number of such sequences is
\[
n^k.
\]

This model is used when:
\begin{itemize}
    \item order is recorded,
    \item each position is chosen separately,
    \item repetition is allowed,
    \item the outcome is a full sequence.
\end{itemize}

Typical examples include PIN codes, repeated coin tosses, repeated die rolls, strings over an alphabet, and repeated selections with replacement.

\begin{examplebox}
The sample space for three tosses of a coin is
\[
\Omega=\{H,T\}^3.
\]
It has
\[
2^3=8
\]
elementary outcomes. Each outcome is a sequence, for example
\[
(H,T,H).
\]
The outcomes \((H,T,H)\) and \((T,H,H)\) are different because the order of tosses is recorded.
\end{examplebox}

\section{Ordered selections without repetition}

Sometimes only some objects are chosen, but the order or position of the chosen objects matters.

\begin{definitionbox}
An \textbf{ordered selection without repetition} of length \(k\) from \(n\) distinct objects is an ordered list of \(k\) different objects chosen from the \(n\) available objects.
\end{definitionbox}

The number of ordered selections without repetition is
\[
P(n,k)=n(n-1)\cdots(n-k+1)=\frac{n!}{(n-k)!}.
\]

This model is used when:
\begin{itemize}
    \item only some objects are selected,
    \item order or position matters,
    \item no object may be used twice.
\end{itemize}

\begin{examplebox}
Gold, silver, and bronze medals are assigned among \(12\) runners. An outcome is not a set of three runners. It is an ordered triple
\[
(\text{gold},\text{silver},\text{bronze}).
\]
There are
\[
12\cdot 11\cdot 10=P(12,3)
\]
possible assignments.
\end{examplebox}

In some European terminology, ordered selections are called variations. In these notes we use the name ordered selections without repetition, or \(k\)-permutations, because it describes directly what the outcome is.

\section{Arrangements of all objects}

A permutation is an arrangement of all objects in a row.

\begin{definitionbox}
A \textbf{permutation} of \(n\) distinct objects is a linear arrangement of all \(n\) objects.
\end{definitionbox}

There are
\[
n!
\]
permutations of \(n\) distinct objects.

This is a special case of ordered selection without repetition with \(k=n\):
\[
P(n,n)=n!.
\]

\begin{examplebox}
If \(7\) students are arranged in a line, then one outcome is a complete ordered list of all students. The number of arrangements is
\[
7!.
\]
\end{examplebox}

Permutations are used when:
\begin{itemize}
    \item all objects are used,
    \item order matters,
    \item the objects are distinguishable,
    \item the arrangement is linear.
\end{itemize}

\section{Circular arrangements}

A circular arrangement is different from a linear arrangement because rotating the whole arrangement does not create a new relative seating order.

\begin{examplebox}
Suppose four people \(A,B,C,D\) sit around a round table. The linear lists
\[
(A,B,C,D),\qquad (B,C,D,A),\qquad (C,D,A,B),\qquad (D,A,B,C)
\]
represent the same circular arrangement if only relative positions matter.
\end{examplebox}

For \(n\) distinct objects placed around a circle, with rotations identified, the number of circular arrangements is
\[
(n-1)!.
\]

One way to see this is to fix one object as a reference point. The remaining \(n-1\) objects can then be arranged around it in
\[
(n-1)!
\]
ways.

\begin{remarkbox}
A circular arrangement requires a modelling decision. If seats are numbered or if there is a distinguished position, then rotations may become different outcomes and the linear count may be appropriate.
\end{remarkbox}
